% --- IMPORTS ---
\documentclass[leqno]{article}
\usepackage{graphicx}
\usepackage{dsfont}
\usepackage{mathtools}
\usepackage{amssymb}
\usepackage{amsmath}
\usepackage{amsthm}
\usepackage{float}
\usepackage{ragged2e}
\usepackage{tensor}
\usepackage{fancyhdr}
\usepackage{empheq}
\usepackage[most]{tcolorbox}
\usepackage{enumitem}
\usepackage{dirtytalk}
\usepackage{marginnote}
\usepackage{microtype}
\usepackage[
  a4paper,
  left=0.8in,
  right=1in,
  top=1in,
  bottom=0.5in,
  headheight=14pt,
  headsep=0.4in
]{geometry}
\setlength{\parindent}{0cm}
% --- TikZ Packages ---
\usepackage{pgfplots}
\pgfplotsset{compat=1.18}
\usepgfplotslibrary{fillbetween, polar}
\usetikzlibrary{calc, positioning}
\usetikzlibrary{angles, quotes}
\usetikzlibrary{arrows.meta}
\usetikzlibrary{decorations.pathreplacing, decorations.markings}
\usetikzlibrary{patterns}
\usetikzlibrary{intersections}
% --- URL Packages ---
\usepackage[colorlinks=true,linkcolor=red,citecolor=magenta,urlcolor=black]{hyperref}%
\urlstyle{same}
% --- END OF IMPORTS ---

% --- CUSTOM COMMANDS ---
\RenewDocumentCommand{\marks}{o m}{%
  \IfValueTF{#1}%
    {\marginnote{\raggedright\textbf{[#2]}}[#1]}%
    {\marginnote{\raggedright\textbf{[#2]}}}%
}
\newcommand{\vspacerule}[2]{%
  \par\vspace*{#1}%
  \noindent\hrulefill\par
  \vspace*{#2}%
}
\definecolor{scarlet}{RGB}{205, 40, 75}
\newcommand{\questionitem}{%
  \item\phantomsection
  \addcontentsline{toc}{section}{Question \number\value{enumi}}%
}
\renewcommand{\contentsname}{Questions}
% --- END OF CUSTOM COMMANDS ---

% --- HEADER CONFIG ---
\pagestyle{fancy}
\fancyhead{}
\fancyfoot{}
\renewcommand{\headrulewidth}{0.9pt}
\lhead{Hyperbolic Functions}
\chead{}
\rhead{\href{https://danieljsmith.org}{danieljsmith.org}}
% --- END OF HEADER CONFIG ---

\begin{document}
\vspace*{20pt}
\tableofcontents
\newpage
\begin{enumerate}[
  label=\textbf{\arabic*.},
  leftmargin=*,
  topsep=0pt,
  partopsep=0pt,
  parsep=0pt
]

% ---- Question 1 ----
\questionitem
% [Source: _QBT__Hyperbolic_Functions, Q1]

The curve $C_1$ has equation \[y=2\cosh 2x-2\] and the curve $C_2$ has equation \[y=3-5e^{-2x}\] \\
\begin{enumerate}[label=\textbf{(\alph*)}]
  \item Sketch the graphs of $C_1$ and $C_2$ on one set of axes, giving the equation of any asymptote and the coordinates of the points where the curves meet the axes. \marks{4} \\
  \item Solve the equation $2\cosh 2x-2=3-5e^{-2x}$, giving your answers in the form $\frac{1}{2}\ln k$, where $k$ is an integer. \marks{5} \\
\end{enumerate}

\newpage

% ---- Question 2 ----
\questionitem
% [Source: _QBT__Hyperbolic_Functions, Q2]

\begin{enumerate}[label=\textbf{(\alph*)}]
  \item Use the definitions of $\sinh x$ and $\cosh x$ in terms of exponentials to show that
  \[
  \tanh x \equiv \frac{e^{2x}-1}{e^{2x}+1}
  \]
  \marks[-20pt]{2}
  \item Hence find the value of $x$ for which
  \[
  \tanh(x-\ln 2)=\frac{3}{7}
  \]
  giving your answer in the form $\dfrac{1}{2}\ln k$, where $k$ is a rational number to be determined. \marks{5} \\
\end{enumerate}

\newpage

% ---- Question 3 ----
\questionitem
% [Source: _QBT__Hyperbolic_Functions, Q3]

\begin{enumerate}[label=\textbf{(\alph*)}]
  \item Express $7\cosh x - 5\sinh x$ in the form $R\cosh(x - \alpha)$, where $R > 0$ and $\alpha > 0$.\\ Give $\alpha$ to $3$ decimal places. \marks{4} \\
  \item Hence solve the equation $7\cosh x - 5\sinh x = 10$, giving all solutions correct to $2$ decimal places. \marks{3} \\
  \item Solve $7\cosh x - 5\sinh x = 10$ by instead using the exponential definitions of $\cosh x$ and $\sinh x$. \marks{4} \\
\end{enumerate}

\newpage

% ---- Question 4 ----
\questionitem
% [Source: _QBT__Hyperbolic_Functions, Q4]

\begin{enumerate}[label=\textbf{(\alph*)}]
  \item Starting from the identity $\cosh 2x = 2\cosh^2 x - 1$, prove that
  \[
  \operatorname{sech} 2x = \frac{\operatorname{sech}^2 x}{2-\operatorname{sech}^2 x}
  \]
  \marks[-30pt]{2}
  \item The function $f$ is defined by
  \[
  f(x)=\operatorname{sech} x \qquad (x>0)
  \]
  \begin{enumerate}[label=(\roman*)]
    \item State the range of $f$. \marks{1}
    \item Use part (a) and part (b)(i) to prove that $\operatorname{sech} 2x<\operatorname{sech} x$ if $x>0$. \marks{3} \\
  \end{enumerate}
\end{enumerate}

\newpage

% ---- Question 5 ----
\questionitem
% [Source: _QBT__Hyperbolic_Functions, Q5]

By writing $\sinh x$ and $\cosh x$ in terms of exponentials, solve the equation \[2\sinh x+5\cosh x=6\] giving your solutions in exact logarithmic form. \marks{6}

\newpage


% ---- Question 6 ----
\questionitem
% [Source: _QBT__Hyperbolic_Functions, Q6]

\begin{enumerate}[label=\textbf{(\alph*)}]
  \item Using the definitions of $\sinh x$ and $\cosh x$ in terms of exponentials, show that \[\cosh 2x = 1+2\sinh^2 x\] \marks[-20pt]{2} 
  \item Hence solve the equation \[2\cosh 2x-7\sinh x=4\] giving all your answers in logarithmic form. \marks{5} \\
\end{enumerate}

\newpage

% ---- Question 7 ----
\questionitem
% [Source: _QBT__Hyperbolic_Functions, Q7]

\begin{enumerate}[label=\textbf{(\alph*)}]
  \item Prove that, for $x\ge 1$,
  \[
  \operatorname{arcosh} x = \ln\left(x+\sqrt{x^2-1}\right)
  \]
  \marks[-20pt]{5}
  \item Prove that the graphs of
  \[
  y=\operatorname{arsinh} x \quad \text{and} \quad y=\operatorname{arcosh}(x+2)
  \]
  do not intersect. \marks{3} \\
\end{enumerate}

\newpage

% ---- Question 8 ----
\questionitem
% [Source: _QBT__Hyperbolic_Functions, Q8]

The function $f(x)$ is defined by $f(x)=\ln(\cosh x-2\sinh x)$. \\
\begin{enumerate}[label=\textbf{(\alph*)}]
  \item Given that $k$ lies in the domain of this function, explain why $k<\frac{1}{2}\ln 3$. \marks{2} \\
  \item
  \begin{enumerate}[label=(\roman*)]
    \item Find $f'(x)$. \marks{2}
    \item Show that
    \[
    f''(x)=\frac{a}{(\cosh x-2\sinh x)^2}
    \]
    where $a$ is an integer to be determined. \marks{3} \\
  \end{enumerate}
  \item Hence find a quadratic approximation to $f(x)$ for small values of $x$. \marks{3} \\
  \item Find the percentage error in this approximation when $x=-0.1$. \marks{2} \\
\end{enumerate}

\newpage

% ---- Question 9 ----
\questionitem
% [Source: _QBT__Hyperbolic_Functions, Q9]

\begin{enumerate}[label=\textbf{(\alph*)}]
  \item Starting from the definitions of $\sinh x$ and $\cosh x$ in terms of exponentials, prove that
  \[
  \cosh x + \sinh x = e^x \quad \text{and} \quad \cosh x - \sinh x = e^{-x}
  \]
  \marks[-20pt]{3}
  \item Solve the equation
  \[
  4\cosh x - 3\sinh x = 5
  \]
  giving your answers as exact logarithms. \marks{5} \\
\end{enumerate}

\newpage

% ---- Question 10 ----
\questionitem
% [Source: _QBT__Hyperbolic_Functions, Q10]

Using definitions of the hyperbolic functions in terms of exponentials, prove that, for all real numbers $x$ and $y$,
\[
\tanh(x+y)=\frac{\tanh x+\tanh y}{1+\tanh x\tanh y}
\]
\marks[-30pt]{5}

\newpage

% ---- Question 11 ----
\questionitem
% [Source: _QBT__Hyperbolic_Functions, Q11]

\begin{enumerate}[label=\textbf{(\alph*)}]
  \item Using the definition of $\tanh x$ in terms of exponentials, prove that
  \[
  \tanh 3x \equiv \frac{3\tanh x + \tanh^3 x}{1 + 3\tanh^2 x}
  \]
  \marks[-30pt]{3}
  \item Hence solve the equation
  \[
  \tanh 3x = 2\tanh x
  \]
  giving your answers as simplified natural logarithms where appropriate. \marks{4} \\
\end{enumerate}

\newpage

% ---- Question 12 ----
\questionitem
% [Source: _QBT__Hyperbolic_Functions, Q12]

\begin{figure}[H]
    \centering
\begin{tikzpicture}[scale=1.2]
\def\tmin{-1.2}
\def\tmax{1.2}
\def\xmin{-2.5}
\def\xmax{4.8}
\def\ymin{-0.4}
\def\ymax{7.5}
\def\tB{ln(2)}

\draw[-{Stealth[length=3mm]}, thick] (\xmin,0) -- (\xmax,0) node[right] {$x$};
\draw[-{Stealth[length=3mm]}, thick] (0,\ymin) -- (0,\ymax) node[above] {$y$};

\draw[thick, black, name path=curve, samples=200, domain=\tmin:\tmax, variable=\t]
  plot ({1 + 2*sinh(\t)},{3*cosh(\t) + sinh(\t)});

\path[name path=yaxis] (0,\ymin) -- (0,\ymax);
\path[name intersections={of=curve and yaxis, by=A}];

\coordinate (B) at ({1 + 2*sinh(\tB)},{3*cosh(\tB) + sinh(\tB)});

\node[below left] at (0,0) {$O$};
\node[below left] at (A) {$A$};
\node[right] at (B) {$B$};
\end{tikzpicture}
\end{figure}


The diagram shows the curve with parametric equations
\[
x = 1 + 2\sinh t, \quad y = 3\cosh t + \sinh t
\] 
\begin{enumerate}[label=\textbf{(\alph*)}]
  \item The curve crosses the positive $y$-axis at $A$.
  \begin{enumerate}[label=(\roman*)]
    \item Determine the value of the parameter $t$ at $A$, giving your answer in logarithmic form. \marks{4}
    \item Find the $y$-coordinate of $A$, giving your answer correct to $3$ significant figures. \marks{2} \\
  \end{enumerate}
  \item The point $B$ has parameter $t = \ln 2$. \\

  Determine the equation of the tangent to the curve at $B$. \marks{6} \\
\end{enumerate}

\newpage

% ---- Question 13 ----
\questionitem
% [Source: _QBT__Hyperbolic_Functions, Q13]

\begin{enumerate}[label=\textbf{(\alph*)}]
  \item It is given that, for $|t|<1$,
  \[
  q = \frac{1}{2}\ln\left(\frac{1+t}{1-t}\right)
  \]
  Starting from the exponential definition of the $\tanh$ function, show that
  \[
  \tanh q = t
  \]
  \marks[-20pt]{4}
  \item Solve the equation
  \[
  6\operatorname{sech}^2 x = 5 - \tanh x
  \]
  Give your answers in logarithmic form. \marks{4} \\
\end{enumerate}

\newpage

% ---- Question 14 ----
\questionitem
% [Source: _QBT__Hyperbolic_Functions, Q14]

\begin{enumerate}[label=\textbf{(\alph*)}]
  \item Sketch the graph of $y = \sinh x$. \marks{2} \\
  \item Given that $u = \sinh x$, use the definition of $\sinh x$ in terms of $e^x$ and $e^{-x}$ to show that
  \[
  x = \ln\left(u + \sqrt{u^2+1}\right)
  \]
  \marks[-30pt]{6}
  \item
  \begin{enumerate}[label=(\roman*)]
    \item Show that the equation
    \[
    9\cosh^2 x - 24\sinh x = -7
    \]
    can be written as
    \[
    9\sinh^2 x - 24\sinh x + 16 = 0
    \]
    \marks[-20pt]{2}
    \item Hence show that the equation has only one real solution for $x$. \\

    Find this solution in the form $\ln a$, where $a$ is an integer. \marks{5} \\
  \end{enumerate}
\end{enumerate}

\newpage

% ---- Question 15 ----
\questionitem
% [Source: _QBT__Hyperbolic_Functions, Q15]

\begin{enumerate}[label=\textbf{(\alph*)}]
  \item Sketch the graph of $y=\operatorname{cosech} x$ and state the equations of its asymptotes. \marks{3} \\
  \item Use the definitions of $\sinh x$ and $\cosh x$ in terms of $e^x$ and $e^{-x}$ to show that, for $x\ne 0$,
  \[
  \coth^2 x - \operatorname{cosec h}^2 x = 1
  \]
  \marks[-20pt]{3}
  \item Solve the equation $\operatorname{cosech}^2 x = 5 - \coth x$, giving your answers in terms of natural logarithms. \marks{5} \\
\end{enumerate}

\newpage

% ---- Question 16 ----
\questionitem
% [Source: _QBT__Hyperbolic_Functions, Q16]

The hyperbola $H$ has equation
\[
\frac{y^2}{b^2}-\frac{x^2}{a^2}=1
\] \\
\begin{enumerate}[label=\textbf{(\alph*)}]
  \item Use calculus to show that the equation of the tangent to $H$ at the point $(a\sinh\theta,\, b\cosh\theta)$ may be written in the form
  \[
  ay\cosh\theta-bx\sinh\theta=ab
  \]
  \marks[-20pt]{4}
\end{enumerate}
The line $\ell_1$ is the tangent to $H$ at the point $(a\sinh\theta,\, b\cosh\theta)$, where $\theta>0$. \\
Given that $\ell_1$ meets the asymptote $y=\frac{b}{a}x$ at the point $P$ \\
\begin{enumerate}[label=\textbf{(\alph*)}, resume]
  \item find, in terms of $a$, $b$ and $\theta$, the coordinates of $P$. \marks{2} \\
\end{enumerate}
The line $\ell_2$ is the tangent to $H$ at the point $(0,\,b)$. \\
Given that $\ell_1$ and $\ell_2$ meet at the point $Q$ \\
\begin{enumerate}[label=\textbf{(\alph*)}, resume]
  \item find, in terms of $a$, $b$ and $\theta$, the coordinates of $Q$. \marks{2} \\
  \item Show that, as $\theta$ varies, the locus of the mid-point of $PQ$ has equation
  \[
  2ay^2-2bxy=ab^2
  \]
  \marks[-20pt]{6}
\end{enumerate}

\newpage

% ---- Question 17 ----
\questionitem
% [Source: _QBT__Hyperbolic_Functions, Q17]

\begin{enumerate}[label=\textbf{(\alph*)}]
  \item Show, by means of a sketch, that the curves with equations
  \[
  y = \tanh x
  \]
  and
  \[
  y = \frac{1}{2}\bigl(1 + \operatorname{sech} x\bigr)
  \]
  have exactly one point of intersection. \marks{4} \\
  \item Find the $x$-coordinate of this point of intersection, giving your answer in the form $a \ln b$. \marks{4} \\
\end{enumerate}
\newpage
% ---- Question 18 ----
\questionitem
% [Source: _QBT__Hyperbolic_Functions, Q18]

\begin{enumerate}[label=(\roman*)]
  \item Sketch the graph of $y=\cosh x$ for $x\ge 0$ and state the value of the gradient when $x=0$. On the same axes, sketch the graph of $y=\operatorname{arcosh} x$ for $x\ge 1$. Label each curve and give the equation of the line of symmetry. \marks{4} \\
  \item Find
  \[
  \int_0^k \cosh x\,\mathrm{d}x
  \]
  where $k>0$. \marks{2} \\
  \item Deduce, or otherwise show, that
  \[
  \int_1^{\cosh k} \operatorname{arcosh} x\,\mathrm{d}x = k\cosh k - \sinh k
  \]
  \marks[-30pt]{4}
\end{enumerate}


\end{enumerate}
\end{document}
